\documentclass{article}
\usepackage{graphicx} % Required for inserting images
\usepackage[T1]{fontenc}
\usepackage[margin=1in]{geometry}
\usepackage{booktabs}
\usepackage{tabularx}
\usepackage[hidelinks]{hyperref}

\title{Position Size and Exit Rules for Single-Leg Delta-Hedged SPY Options}
\author{Strategy research \textperiodcentered{} L/S Gamma Desk, QUANTT}
\date{2026-10-01}

\begin{document}

\maketitle

\textbf{Question:} for the size and exit rules lets also run a /strategy-research on this

\textbf{Method:} Searched all 27 papers in \texttt{docs/\allowbreak{}papers/\allowbreak{}} (page-tagged index) + web search + the desk's v1 rules (\texttt{archive/\allowbreak{}v1-\allowbreak{}ls-\allowbreak{}gamma} README). Quotes verified against the cited PDF pages. Page numbers are PDF pages; printed page numbers are given where they differ.

\textbf{Related research:} \href{2026-10-01-target-dte-moneyness-calls-vs-puts.pdf}{Target DTE, moneyness, calls vs puts} \textperiodcentered{} \href{2026-09-29-signal-generation-for-market-state.pdf}{Signal generation}

\textbf{Desk decisions already made (PM, 2026-10-01):} v1 VRP band; single-leg calls and puts; daily delta hedging.

\section{Short answer}

The papers agree on \textbf{four kinds of exit}:

\begin{enumerate}

\item \textbf{Time:} get out before the final week.

\item \textbf{Signal reversal:} the thesis is gone.

\item \textbf{Loss limit:} a hard risk stop.

\item \textbf{Profit take:} optional.

\end{enumerate}

For \textbf{size}, the serious frameworks either cap a risk measure (vega, max loss, collateral) or use a fraction of Kelly. Kelly needs P\&L statistics the desk doesn't have yet.

For paper trading, the simplest sound start is \textbf{1 contract per trade} with time, signal and loss exits. Scale up only once real P\&L data exists.

\section{What the papers say}

\subsection{Time exits: avoid the final week}

\begin{itemize}

\item \textbf{PRISM rolls before the last week:}

\begin{quote}``VRC rolls all options positions when time to expiry falls below 5 trading days.'' (Verma 2026, \emph{PRISM}, p. 36)\end{quote}

The reason given is that it:

\begin{quote}``avoids the theta cliff of the final week'' (\emph{PRISM}, p. 36)\end{quote}

\item \textbf{v1 desk rule:} enter at 7--28 DTE and exit no later than 3 DTE (v1 README).

\item \textbf{The Cboe PutWrite benchmark is the opposite extreme:} it sells 1-month ATM puts monthly and \textbf{holds them to expiry}, fully collateralized (\href{https://en.wikipedia.org/wiki/CBOE_S\&P_500_PutWrite_Index}{Wikipedia}, \href{https://cdn.cboe.com/api/global/us_indices/governance/Cboe_SP_500_PutWrite_Indices_Methodology.pdf}{Cboe methodology}). That index isn't delta-hedged, though.

\end{itemize}

\subsection{Signal exits: close when the thesis reverses}

\begin{itemize}

\item \textbf{PRISM's VRP reversal trigger} fires when the VRP z-score stays negative for three consecutive closes. This:

\begin{quote}``signals that implied vol has fallen below VRC's regime-adjusted realized vol estimate'' (\emph{PRISM}, p. 37)\end{quote}

\item \textbf{OPHR's baselines close and flip:}

\begin{quote}``If predicted volatility crosses the opposite threshold, the position is closed and reversed.'' (Chen, Cai, Qin \& An, \emph{OPHR}, p. 28)\end{quote}

\end{itemize}

\subsection{Loss and profit limits}

\begin{itemize}

\item \textbf{OPHR's baselines use percentage limits plus a maximum hold:}

\begin{quote}``the position is closed if the profit reaches p\% or if the loss reaches l\%. In addition, a maximum holding period'' (\emph{OPHR}, p. 26)\end{quote}

Their settings (take-profit 5\%, stop 3\%, 96-hour maximum hold, 60--90 day options, p. 27) are tuned for hourly crypto trading and \textbf{don't transfer} to daily SPY.

\item \textbf{v1 desk rules:} for long gamma, a VRP hard stop and a realized-vol trailing check; for short gamma, a stop at 30\% of premium.

\item \textbf{Profit taking at 50\% of premium} on short puts is common practitioner advice. I couldn't verify a rigorous study (the one found couldn't be accessed), so treat it as untested.

\end{itemize}

\subsection{Sizing}

\begin{itemize}

\item \textbf{Fractional Kelly:}

\begin{quote}``VRC applies a fractional Kelly factor of 1/2 as standard policy'' (\emph{PRISM}, p. 23)\end{quote}

Kelly needs the mean and variance of trade P\&L, which the desk will only have after paper trading.

\item \textbf{PRISM's hard limits (p. 38):}

\begin{itemize}

\item net vega at most 3\% of assets

\item 99\% daily CVaR at most 0.75\% of notional

\item delta exposure at most 2\ensuremath{\times} assets

\item exit if ATM bid-ask spreads exceed 100bp of vol

\end{itemize}

\item \textbf{The PutWrite index} sizes so that T-bills ``can finance the maximum possible loss from final settlement'' of the puts sold. That's the most conservative rule.

\item \textbf{v1 desk caps:} long gamma at 5--10\% of capital; a 20\% budget for the short side.

\end{itemize}

\section{How the papers connect}

\begin{center}\small
\begin{tabularx}{\linewidth}{>{\raggedright\arraybackslash}X >{\raggedright\arraybackslash}X >{\raggedright\arraybackslash}X}
\toprule
\textbf{Source} & \textbf{Exit approach} & \textbf{Sizing approach} \\
\midrule
PRISM (whitepaper) & Roll at under 5 trading days; signal reversal; P\&L at 2.5\ensuremath{\sigma} below expectation; hard limits & Half-Kelly with an uncertainty penalty; vega, CVaR and leverage caps \\
OPHR baselines & Take-profit, stop-loss, max hold; close and reverse on opposite signal & By margin use \\
Cboe PutWrite (web) & Hold to expiry & Fully collateralized \\
v1 desk & Exit at 3 DTE; VRP and premium stops & 5--10\% long, 20\% short budget \\
\bottomrule
\end{tabularx}
\end{center}

\begin{itemize}

\item \textbf{Everyone uses a time exit and a thesis exit.} They differ mostly on whether to add profit and loss stops.

\item \textbf{Sizing goes from simplest (full collateral) to most data-hungry (Kelly).} Risk caps such as vega or max loss sit in between.

\end{itemize}

\section{Relevance to the LS Gamma desk}

Suggested starting rules for paper trading. These are suggestions; the PM sets the numbers.

\begin{enumerate}

\item \textbf{Size:} 1 contract per trade. Paper trading is first a mechanics test. Later, size by a vega or max-loss budget, and only consider half-Kelly once there's P\&L history.

\item \textbf{Time exit:} close at about 7 calendar DTE, which is roughly PRISM's 5 trading days.

\item \textbf{Signal exit:} close when the signal is back to flat for 3 consecutive days. Close and flip on an opposite signal.

\item \textbf{Loss stop:}

\begin{itemize}

\item short put: when the position loss reaches 1--2\ensuremath{\times} the premium received

\item long call: at 50\% of the premium paid

\end{itemize}

Measure losses including the hedge P\&L.

\item \textbf{Optional profit take:} close short puts at 50\% of premium. Untested; consider running it in shadow to compare.

\item \textbf{Log every exit and its trigger,} so the rules can be tuned from data.

\end{enumerate}

\section{Gaps and open questions}

\begin{itemize}

\item \textbf{No paper tests these exits on daily-hedged single SPY options.} The OPHR numbers are crypto; PRISM is simulated; PutWrite is unhedged.

\item \textbf{Stops and daily hedging interact.} A stop on option P\&L alone ignores hedge gains, so P\&L should be measured on the combined position.

\item \textbf{The account's starting capital and risk budget} aren't set yet, and they're needed before scaling past 1 contract.

\end{itemize}

\section{Sources}

\textbf{Repo papers used, with pages cited:}

\begin{itemize}

\item \texttt{PRISM-\allowbreak{}A-\allowbreak{}Probabilistic-\allowbreak{}Regime-\allowbreak{}Integrated-\allowbreak{}Scalping-\allowbreak{}Model-\allowbreak{}for-\allowbreak{}Derivatives-\allowbreak{}Arbitrage.\allowbreak{}pdf} (pp. 23, 36, 37, 38)

\item \texttt{OPHR-\allowbreak{}Mastering-\allowbreak{}Volatility-\allowbreak{}Trading-\allowbreak{}with-\allowbreak{}Multi-\allowbreak{}Agent-\allowbreak{}Deep-\allowbreak{}Reinforcement-\allowbreak{}Learning.\allowbreak{}pdf} (pp. 26, 27, 28)

\end{itemize}

\textbf{Checked, nothing on sizing or exits:} the deep hedging papers (they size hedges, not positions), the forecasting papers, \texttt{Agents-\allowbreak{}Are-\allowbreak{}not-\allowbreak{}Algorithms.\allowbreak{}pdf}, \texttt{From-\allowbreak{}Deep-\allowbreak{}Learning-\allowbreak{}to-\allowbreak{}LLMs-\allowbreak{}A-\allowbreak{}Survey-\allowbreak{}of-\allowbreak{}AI-\allowbreak{}in-\allowbreak{}Quantitative-\allowbreak{}Investment.\allowbreak{}pdf}, Sepp (hedging frequency, not exits).

\textbf{Web sources (outside the repo):}

\begin{itemize}

\item \href{https://cdn.cboe.com/api/global/us_indices/governance/Cboe_SP_500_PutWrite_Indices_Methodology.pdf}{Cboe S\&P 500 PutWrite Indices methodology}

\item \href{https://en.wikipedia.org/wiki/CBOE_S\&P_500_PutWrite_Index}{CBOE S\&P 500 PutWrite Index (Wikipedia)}

\item \href{https://cdn.cboe.com/resources/education/research_publications/PUTIndexEnnisKnupp.pdf}{Ennis Knupp, \emph{Evaluating the Performance Characteristics of the CBOE S\&P 500 PutWrite Index}}

\end{itemize}

\textbf{Suggested papers to add to \texttt{docs/\allowbreak{}papers/\allowbreak{}}:}

\begin{enumerate}

\item \textbf{Ennis Knupp PUT index evaluation:} a long-run benchmark for a systematic ATM put-selling program.

\item \textbf{A rigorous study of profit and loss rules for short index options,} if one can be found. The practitioner claims remain unverified.

\end{enumerate}

\end{document}
